% --- Archon physics formalization source begin ---
% archon:physics
% archon:covers PhyXMiniProblems/problem_phyx_mini_0491.lean
% archon:source-report reports/phyx_mini/problem_phyx_mini_0491.source.json

\chapter{Physics problem phyx\_mini\_0491}
\label{ch:PhyXMiniProblems_problem_phyx_mini_0491}

\paragraph{Problem source.}
\#\# Physical scenario

An ideal gas is slowly compressed at a constant pressure of \textbackslash{}(2.0\textbackslash{},\textbackslash{}mathrm\{atm\}\textbackslash{}) from \textbackslash{}(10.0\textbackslash{},\textbackslash{}mathrm\{L\}\textbackslash{}) to \textbackslash{}(2.0\textbackslash{},\textbackslash{}mathrm\{L\}\textbackslash{}). This process is represented in figure as the path B to D. (In this process, some heat flows out of the gas and the temperature drops.) Heat is then added to the gas, holding the volume constant, and the pressure and temperature are allowed to rise (line DA) until the temperature reaches its original value \textbackslash{}((T\_A = T\_B)\textbackslash{}).

\#\# Question

In the process BDA, calculate the total heat flow into the gas.

\#\# Answer choices

- A: A: -1800J

- B: B: 1800J

- C: C: 1600J

- D: D: -1600J

\#\# Figure caption (auxiliary; use the image as primary evidence)

The image is a pressure-volume (P-V) diagram illustrating a thermodynamic cycle. It features four distinct points labeled A, B, D, and a process labeled between specific points. The axes are labeled "P" for pressure and "V (L)" for volume in liters.

1. **Points and Lines:**
   - **A** is at the top left, connected horizontally to **D**.
   - **B** is at the bottom right, connected vertically to **D**.
   - A curve connects **A** to **B**, labeled "Isothermal."
   - A horizontal line labeled "Isobaric" connects **D** to **B**.
   - A vertical line labeled "Isovolumetric" connects **D** to **A**.

2. **Text and Labels:**
   - The vertical axis indicates pressure with values, and key pressures are labeled as \textbackslash{}( P\_A \textbackslash{}) and \textbackslash{}( P\_B \textbackslash{}).
   - The horizontal axis indicates volume in liters, ranging from 0 to 10.
   - Labels for processes: "Isovolumetric" near the vertical line, "Isobaric" near the horizontal line, and "Isothermal" along the curve.

3. **Arrows:**
   - An upward arrow next to the "Isovolumetric" line and a left-pointing arrow along the "Isobaric" line indicate the direction of the processes.

The diagram represents a cycle involving isochoric (constant volume), isobaric (constant pressure), and isothermal (constant temperature) processes.

\#\# Dataset metadata

Laws of Thermodynamics; Multi-Formula Reasoning, Physical Model Grounding Reasoning

\paragraph{Recorded answer/context.}
D: D: -1600J

\paragraph{Figure/image path.}
/root/proposal\_for\_physic/science-mango/phyx\_mini\_run/phyx\_data/test\_image/491.png

\paragraph{Formalization target.}
create a compiling Lean file with sorry bodies at `PhyXMiniProblems/problem\_phyx\_mini\_0491.lean`. The Lean declarations must preserve the physical quantities, dimensions or dimensional roles, figure labels, governing-law hypotheses, and final relation expressed by this problem.
Use Mathlib/Physlib names found through LeanExplore where available. If a physics API is missing, introduce faithful local abstractions rather than scalar placeholder aliases.

\begin{theorem}[Physics formalization target]
\label{thm:physics:phyx_mini_0491:target}
\lean{PhyXMiniProblems.ProblemPhyXMini0491.problem_phyx_mini_0491}
\uses{def:physics:phyx-mini-0491:phyxminiproblems-problemphyxmini0491-idealgasbdaprocess, def:physics:phyx-mini-0491:phyxminiproblems-problemphyxmini0491-totalheatintogasinjoules, def:physics:phyx-mini-0491:phyxminiproblems-problemphyxmini0491-matchesproblemandprimaryfigure, def:physics:phyx-mini-0491:phyxminiproblems-problemphyxmini0491-hasphysicalparameters, def:physics:phyx-mini-0491:phyxminiproblems-problemphyxmini0491-satisfiesidealgasprocesslaws, def:physics:phyx-mini-0491:phyxminiproblems-problemphyxmini0491-displayedtotalheatinjoules, def:physics:phyx-mini-0491:phyxminiproblems-problemphyxmini0491-recordedanswerchoice, def:physics:phyx-mini-0491:phyxminiproblems-problemphyxmini0491-roundstonearesthundredjoules, def:physics:phyx-mini-0491:phyxminiproblems-problemphyxmini0491-isuniqueclosestdisplayedheat}
The assigned autoformalize agent should translate this physics problem into Lean declarations in the covered file, with theorem and lemma proof bodies written as `by sorry`.
\end{theorem}
\begin{proof}
This is an autoformalization task, not a proof task. Produce statements that can later be proved by the physics prover without weakening the physical model.
\end{proof}

% --- Archon declaration topology begin ---
\section{Declaration topology}
\begin{definition}[Volume Quantity]
  \label{def:physics:phyx-mini-0491:phyxminiproblems-problemphyxmini0491-volumequantity}
  \lean{PhyXMiniProblems.ProblemPhyXMini0491.VolumeQuantity}
  A nonnegative physical volume carrying dimension L³
\end{definition}
\begin{proof}
  This is the declared model definition of Volume Quantity.
\end{proof}
\begin{definition}[volume In Cubic Meters]
  \label{def:physics:phyx-mini-0491:phyxminiproblems-problemphyxmini0491-volumeincubicmeters}
  \lean{PhyXMiniProblems.ProblemPhyXMini0491.volumeInCubicMeters}
  \uses{def:physics:phyx-mini-0491:phyxminiproblems-problemphyxmini0491-volumequantity}
  Cubic-metre readout of a physical volume
\end{definition}
\begin{proof}
  This is the declared model definition of volume In Cubic Meters.
\end{proof}
\begin{definition}[volume In Liters]
  \label{def:physics:phyx-mini-0491:phyxminiproblems-problemphyxmini0491-volumeinliters}
  \lean{PhyXMiniProblems.ProblemPhyXMini0491.volumeInLiters}
  \uses{def:physics:phyx-mini-0491:phyxminiproblems-problemphyxmini0491-volumequantity, def:physics:phyx-mini-0491:phyxminiproblems-problemphyxmini0491-volumeincubicmeters}
  Litre readout of a physical volume, using 1 m³ = 1000 L
\end{definition}
\begin{proof}
  This is the declared model definition of volume In Liters.
\end{proof}
\begin{definition}[pressure In Pascals]
  \label{def:physics:phyx-mini-0491:phyxminiproblems-problemphyxmini0491-pressureinpascals}
  \lean{PhyXMiniProblems.ProblemPhyXMini0491.pressureInPascals}
  Pascal readout of a dimensionful pressure
\end{definition}
\begin{proof}
  This is the declared model definition of pressure In Pascals.
\end{proof}
\begin{definition}[pressure In Atmospheres]
  \label{def:physics:phyx-mini-0491:phyxminiproblems-problemphyxmini0491-pressureinatmospheres}
  \lean{PhyXMiniProblems.ProblemPhyXMini0491.pressureInAtmospheres}
  Standard-atmosphere readout of a dimensionful pressure
\end{definition}
\begin{proof}
  This is the declared model definition of pressure In Atmospheres.
\end{proof}
\begin{definition}[energy In Joules]
  \label{def:physics:phyx-mini-0491:phyxminiproblems-problemphyxmini0491-energyinjoules}
  \lean{PhyXMiniProblems.ProblemPhyXMini0491.energyInJoules}
  Joule readout of a signed physical energy, heat transfer, or work
\end{definition}
\begin{proof}
  This is the declared model definition of energy In Joules.
\end{proof}
\begin{definition}[joules Per Liter Atmosphere]
  \label{def:physics:phyx-mini-0491:phyxminiproblems-problemphyxmini0491-joulesperliteratmosphere}
  \lean{PhyXMiniProblems.ProblemPhyXMini0491.joulesPerLiterAtmosphere}
  Exact SI value of one litre-atmosphere, in joules
\end{definition}
\begin{proof}
  This is the declared model definition of joules Per Liter Atmosphere.
\end{proof}
\begin{lemma}[standard Atmosphere In Pascals]
  \label{lem:physics:phyx-mini-0491:phyxminiproblems-problemphyxmini0491-standardatmosphereinpascals}
  \lean{PhyXMiniProblems.ProblemPhyXMini0491.standardAtmosphereInPascals}
  \uses{def:physics:phyx-mini-0491:phyxminiproblems-problemphyxmini0491-pressureinpascals}
  Physlib's standard atmosphere is exactly 101325 Pa
\end{lemma}
\begin{proof}
  The conclusion follows from the governing relations and figure readouts named in the statement of standard Atmosphere In Pascals.
\end{proof}
\begin{definition}[Thermodynamic State]
  \label{def:physics:phyx-mini-0491:phyxminiproblems-problemphyxmini0491-thermodynamicstate}
  \lean{PhyXMiniProblems.ProblemPhyXMini0491.ThermodynamicState}
  The three labeled equilibrium states in the supplied pV diagram
\end{definition}
\begin{proof}
  This is the declared model definition of Thermodynamic State.
\end{proof}
\begin{definition}[Process Leg]
  \label{def:physics:phyx-mini-0491:phyxminiproblems-problemphyxmini0491-processleg}
  \lean{PhyXMiniProblems.ProblemPhyXMini0491.ProcessLeg}
  The two directed legs of the requested process B → D → A
\end{definition}
\begin{proof}
  This is the declared model definition of Process Leg.
\end{proof}
\begin{definition}[leg Source]
  \label{def:physics:phyx-mini-0491:phyxminiproblems-problemphyxmini0491-legsource}
  \lean{PhyXMiniProblems.ProblemPhyXMini0491.legSource}
  \uses{def:physics:phyx-mini-0491:phyxminiproblems-problemphyxmini0491-thermodynamicstate, def:physics:phyx-mini-0491:phyxminiproblems-problemphyxmini0491-processleg}
  Initial state of a directed process leg
\end{definition}
\begin{proof}
  This is the declared model definition of leg Source.
\end{proof}
\begin{definition}[leg Target]
  \label{def:physics:phyx-mini-0491:phyxminiproblems-problemphyxmini0491-legtarget}
  \lean{PhyXMiniProblems.ProblemPhyXMini0491.legTarget}
  \uses{def:physics:phyx-mini-0491:phyxminiproblems-problemphyxmini0491-thermodynamicstate, def:physics:phyx-mini-0491:phyxminiproblems-problemphyxmini0491-processleg}
  Final state of a directed process leg
\end{definition}
\begin{proof}
  This is the declared model definition of leg Target.
\end{proof}
\begin{definition}[Diagram Segment]
  \label{def:physics:phyx-mini-0491:phyxminiproblems-problemphyxmini0491-diagramsegment}
  \lean{PhyXMiniProblems.ProblemPhyXMini0491.DiagramSegment}
  The three red segments or curves visible in the primary raster
\end{definition}
\begin{proof}
  This is the declared model definition of Diagram Segment.
\end{proof}
\begin{definition}[Process Kind]
  \label{def:physics:phyx-mini-0491:phyxminiproblems-problemphyxmini0491-processkind}
  \lean{PhyXMiniProblems.ProblemPhyXMini0491.ProcessKind}
  Thermodynamic constraint named beside a path in the diagram
\end{definition}
\begin{proof}
  This is the declared model definition of Process Kind.
\end{proof}
\begin{definition}[displayed Process Kind]
  \label{def:physics:phyx-mini-0491:phyxminiproblems-problemphyxmini0491-displayedprocesskind}
  \lean{PhyXMiniProblems.ProblemPhyXMini0491.displayedProcessKind}
  \uses{def:physics:phyx-mini-0491:phyxminiproblems-problemphyxmini0491-diagramsegment, def:physics:phyx-mini-0491:phyxminiproblems-problemphyxmini0491-processkind}
  Process label attached to each visible segment
\end{definition}
\begin{proof}
  This is the declared model definition of displayed Process Kind.
\end{proof}
\begin{definition}[leg Process Kind]
  \label{def:physics:phyx-mini-0491:phyxminiproblems-problemphyxmini0491-legprocesskind}
  \lean{PhyXMiniProblems.ProblemPhyXMini0491.legProcessKind}
  \uses{def:physics:phyx-mini-0491:phyxminiproblems-problemphyxmini0491-processleg, def:physics:phyx-mini-0491:phyxminiproblems-problemphyxmini0491-processkind}
  Process kind of each traversed leg
\end{definition}
\begin{proof}
  This is the declared model definition of leg Process Kind.
\end{proof}
\begin{definition}[Process Regime]
  \label{def:physics:phyx-mini-0491:phyxminiproblems-problemphyxmini0491-processregime}
  \lean{PhyXMiniProblems.ProblemPhyXMini0491.ProcessRegime}
  Whether a leg is traversed slowly enough to be modeled quasistatically
\end{definition}
\begin{proof}
  This is the declared model definition of Process Regime.
\end{proof}
\begin{definition}[Segment Shape]
  \label{def:physics:phyx-mini-0491:phyxminiproblems-problemphyxmini0491-segmentshape}
  \lean{PhyXMiniProblems.ProblemPhyXMini0491.SegmentShape}
  Geometric shape of a visible path in the pV plane
\end{definition}
\begin{proof}
  This is the declared model definition of Segment Shape.
\end{proof}
\begin{definition}[displayed Segment Shape]
  \label{def:physics:phyx-mini-0491:phyxminiproblems-problemphyxmini0491-displayedsegmentshape}
  \lean{PhyXMiniProblems.ProblemPhyXMini0491.displayedSegmentShape}
  \uses{def:physics:phyx-mini-0491:phyxminiproblems-problemphyxmini0491-diagramsegment, def:physics:phyx-mini-0491:phyxminiproblems-problemphyxmini0491-segmentshape}
  Shape read directly from the primary raster
\end{definition}
\begin{proof}
  This is the declared model definition of displayed Segment Shape.
\end{proof}
\begin{definition}[Arrow Direction]
  \label{def:physics:phyx-mini-0491:phyxminiproblems-problemphyxmini0491-arrowdirection}
  \lean{PhyXMiniProblems.ProblemPhyXMini0491.ArrowDirection}
  Direction of a visible process arrow
\end{definition}
\begin{proof}
  This is the declared model definition of Arrow Direction.
\end{proof}
\begin{definition}[displayed Arrow Direction]
  \label{def:physics:phyx-mini-0491:phyxminiproblems-problemphyxmini0491-displayedarrowdirection}
  \lean{PhyXMiniProblems.ProblemPhyXMini0491.displayedArrowDirection}
  \uses{def:physics:phyx-mini-0491:phyxminiproblems-problemphyxmini0491-processleg, def:physics:phyx-mini-0491:phyxminiproblems-problemphyxmini0491-arrowdirection}
  Arrow direction on each of the two traversed legs
\end{definition}
\begin{proof}
  This is the declared model definition of displayed Arrow Direction.
\end{proof}
\begin{definition}[Figure Axis]
  \label{def:physics:phyx-mini-0491:phyxminiproblems-problemphyxmini0491-figureaxis}
  \lean{PhyXMiniProblems.ProblemPhyXMini0491.FigureAxis}
  The two Cartesian axes shown in the raster
\end{definition}
\begin{proof}
  This is the declared model definition of Figure Axis.
\end{proof}
\begin{definition}[Axis Quantity]
  \label{def:physics:phyx-mini-0491:phyxminiproblems-problemphyxmini0491-axisquantity}
  \lean{PhyXMiniProblems.ProblemPhyXMini0491.AxisQuantity}
  Physical quantity represented by an axis
\end{definition}
\begin{proof}
  This is the declared model definition of Axis Quantity.
\end{proof}
\begin{definition}[Axis Display Unit]
  \label{def:physics:phyx-mini-0491:phyxminiproblems-problemphyxmini0491-axisdisplayunit}
  \lean{PhyXMiniProblems.ProblemPhyXMini0491.AxisDisplayUnit}
  Unit information literally printed on an axis
\end{definition}
\begin{proof}
  This is the declared model definition of Axis Display Unit.
\end{proof}
\begin{definition}[Pressure Level Label]
  \label{def:physics:phyx-mini-0491:phyxminiproblems-problemphyxmini0491-pressurelevellabel}
  \lean{PhyXMiniProblems.ProblemPhyXMini0491.PressureLevelLabel}
  The two symbolic pressure levels labeled on the vertical axis
\end{definition}
\begin{proof}
  This is the declared model definition of Pressure Level Label.
\end{proof}
\begin{definition}[Pressure Volume Figure]
  \label{def:physics:phyx-mini-0491:phyxminiproblems-problemphyxmini0491-pressurevolumefigure}
  \lean{PhyXMiniProblems.ProblemPhyXMini0491.PressureVolumeFigure}
  \uses{def:physics:phyx-mini-0491:phyxminiproblems-problemphyxmini0491-thermodynamicstate, def:physics:phyx-mini-0491:phyxminiproblems-problemphyxmini0491-processleg, def:physics:phyx-mini-0491:phyxminiproblems-problemphyxmini0491-diagramsegment, def:physics:phyx-mini-0491:phyxminiproblems-problemphyxmini0491-processkind, def:physics:phyx-mini-0491:phyxminiproblems-problemphyxmini0491-segmentshape, def:physics:phyx-mini-0491:phyxminiproblems-problemphyxmini0491-arrowdirection, def:physics:phyx-mini-0491:phyxminiproblems-problemphyxmini0491-figureaxis, def:physics:phyx-mini-0491:phyxminiproblems-problemphyxmini0491-axisquantity, def:physics:phyx-mini-0491:phyxminiproblems-problemphyxmini0491-axisdisplayunit, def:physics:phyx-mini-0491:phyxminiproblems-problemphyxmini0491-pressurelevellabel}
  Literal visible content of the supplied pressure-volume diagram
\end{definition}
\begin{proof}
  This is the declared model definition of Pressure Volume Figure.
\end{proof}
\begin{definition}[Gas Model]
  \label{def:physics:phyx-mini-0491:phyxminiproblems-problemphyxmini0491-gasmodel}
  \lean{PhyXMiniProblems.ProblemPhyXMini0491.GasModel}
  Thermodynamic model stated in the problem
\end{definition}
\begin{proof}
  This is the declared model definition of Gas Model.
\end{proof}
\begin{definition}[Ideal Gas BDAProcess]
  \label{def:physics:phyx-mini-0491:phyxminiproblems-problemphyxmini0491-idealgasbdaprocess}
  \lean{PhyXMiniProblems.ProblemPhyXMini0491.IdealGasBDAProcess}
  \uses{def:physics:phyx-mini-0491:phyxminiproblems-problemphyxmini0491-volumequantity, def:physics:phyx-mini-0491:phyxminiproblems-problemphyxmini0491-thermodynamicstate, def:physics:phyx-mini-0491:phyxminiproblems-problemphyxmini0491-processleg, def:physics:phyx-mini-0491:phyxminiproblems-problemphyxmini0491-processregime, def:physics:phyx-mini-0491:phyxminiproblems-problemphyxmini0491-pressurevolumefigure, def:physics:phyx-mini-0491:phyxminiproblems-problemphyxmini0491-gasmodel}
  The ideal-gas process and its independent observables. Neither heat transfer is defined from an answer choice. workDoneByGasOn is signed positive for expansion, and heatTransferredIntoGasOn is signed positive into the gas
\end{definition}
\begin{proof}
  This is the declared model definition of Ideal Gas BDAProcess.
\end{proof}
\begin{definition}[internal Energy In Joules]
  \label{def:physics:phyx-mini-0491:phyxminiproblems-problemphyxmini0491-internalenergyinjoules}
  \lean{PhyXMiniProblems.ProblemPhyXMini0491.internalEnergyInJoules}
  \uses{def:physics:phyx-mini-0491:phyxminiproblems-problemphyxmini0491-energyinjoules, def:physics:phyx-mini-0491:phyxminiproblems-problemphyxmini0491-thermodynamicstate, def:physics:phyx-mini-0491:phyxminiproblems-problemphyxmini0491-idealgasbdaprocess}
  Joule readout of internal energy at a labeled state
\end{definition}
\begin{proof}
  This is the declared model definition of internal Energy In Joules.
\end{proof}
\begin{definition}[work Done By Gas In Joules]
  \label{def:physics:phyx-mini-0491:phyxminiproblems-problemphyxmini0491-workdonebygasinjoules}
  \lean{PhyXMiniProblems.ProblemPhyXMini0491.workDoneByGasInJoules}
  \uses{def:physics:phyx-mini-0491:phyxminiproblems-problemphyxmini0491-energyinjoules, def:physics:phyx-mini-0491:phyxminiproblems-problemphyxmini0491-processleg, def:physics:phyx-mini-0491:phyxminiproblems-problemphyxmini0491-idealgasbdaprocess}
  Signed joule readout of work done by the gas on a leg
\end{definition}
\begin{proof}
  This is the declared model definition of work Done By Gas In Joules.
\end{proof}
\begin{definition}[heat Transferred Into Gas In Joules]
  \label{def:physics:phyx-mini-0491:phyxminiproblems-problemphyxmini0491-heattransferredintogasinjoules}
  \lean{PhyXMiniProblems.ProblemPhyXMini0491.heatTransferredIntoGasInJoules}
  \uses{def:physics:phyx-mini-0491:phyxminiproblems-problemphyxmini0491-energyinjoules, def:physics:phyx-mini-0491:phyxminiproblems-problemphyxmini0491-processleg, def:physics:phyx-mini-0491:phyxminiproblems-problemphyxmini0491-idealgasbdaprocess}
  Signed joule readout of heat transferred into the gas on a leg
\end{definition}
\begin{proof}
  This is the declared model definition of heat Transferred Into Gas In Joules.
\end{proof}
\begin{definition}[total Heat Into Gas In Joules]
  \label{def:physics:phyx-mini-0491:phyxminiproblems-problemphyxmini0491-totalheatintogasinjoules}
  \lean{PhyXMiniProblems.ProblemPhyXMini0491.totalHeatIntoGasInJoules}
  \uses{def:physics:phyx-mini-0491:phyxminiproblems-problemphyxmini0491-idealgasbdaprocess, def:physics:phyx-mini-0491:phyxminiproblems-problemphyxmini0491-heattransferredintogasinjoules}
  Total heat into the gas over B → D → A, defined only as the sum of the two independent leg heats. This definition contains no numerical answer
\end{definition}
\begin{proof}
  This is the declared model definition of total Heat Into Gas In Joules.
\end{proof}
\begin{definition}[total Work Done By Gas In Joules]
  \label{def:physics:phyx-mini-0491:phyxminiproblems-problemphyxmini0491-totalworkdonebygasinjoules}
  \lean{PhyXMiniProblems.ProblemPhyXMini0491.totalWorkDoneByGasInJoules}
  \uses{def:physics:phyx-mini-0491:phyxminiproblems-problemphyxmini0491-idealgasbdaprocess, def:physics:phyx-mini-0491:phyxminiproblems-problemphyxmini0491-workdonebygasinjoules}
  Total boundary work done by the gas over the two traversed legs
\end{definition}
\begin{proof}
  This is the declared model definition of total Work Done By Gas In Joules.
\end{proof}
\begin{definition}[Matches Problem And Primary Figure]
  \label{def:physics:phyx-mini-0491:phyxminiproblems-problemphyxmini0491-matchesproblemandprimaryfigure}
  \lean{PhyXMiniProblems.ProblemPhyXMini0491.MatchesProblemAndPrimaryFigure}
  \uses{def:physics:phyx-mini-0491:phyxminiproblems-problemphyxmini0491-volumeinliters, def:physics:phyx-mini-0491:phyxminiproblems-problemphyxmini0491-pressureinpascals, def:physics:phyx-mini-0491:phyxminiproblems-problemphyxmini0491-pressureinatmospheres, def:physics:phyx-mini-0491:phyxminiproblems-problemphyxmini0491-displayedsegmentshape, def:physics:phyx-mini-0491:phyxminiproblems-problemphyxmini0491-displayedarrowdirection, def:physics:phyx-mini-0491:phyxminiproblems-problemphyxmini0491-idealgasbdaprocess, def:physics:phyx-mini-0491:phyxminiproblems-problemphyxmini0491-heattransferredintogasinjoules}
  Exact transcription of the problem prose and primary raster. The final temperature equality is endpoint data, not a heat-value assumption
\end{definition}
\begin{proof}
  This is the declared model definition of Matches Problem And Primary Figure.
\end{proof}
\begin{definition}[Has Physical Parameters]
  \label{def:physics:phyx-mini-0491:phyxminiproblems-problemphyxmini0491-hasphysicalparameters}
  \lean{PhyXMiniProblems.ProblemPhyXMini0491.HasPhysicalParameters}
  \uses{def:physics:phyx-mini-0491:phyxminiproblems-problemphyxmini0491-volumeincubicmeters, def:physics:phyx-mini-0491:phyxminiproblems-problemphyxmini0491-pressureinpascals, def:physics:phyx-mini-0491:phyxminiproblems-problemphyxmini0491-idealgasbdaprocess}
  Positivity conditions for the physical equilibrium states
\end{definition}
\begin{proof}
  This is the declared model definition of Has Physical Parameters.
\end{proof}
\begin{definition}[Satisfies Ideal Gas Process Laws]
  \label{def:physics:phyx-mini-0491:phyxminiproblems-problemphyxmini0491-satisfiesidealgasprocesslaws}
  \lean{PhyXMiniProblems.ProblemPhyXMini0491.SatisfiesIdealGasProcessLaws}
  \uses{def:physics:phyx-mini-0491:phyxminiproblems-problemphyxmini0491-volumeincubicmeters, def:physics:phyx-mini-0491:phyxminiproblems-problemphyxmini0491-pressureinpascals, def:physics:phyx-mini-0491:phyxminiproblems-problemphyxmini0491-legsource, def:physics:phyx-mini-0491:phyxminiproblems-problemphyxmini0491-legtarget, def:physics:phyx-mini-0491:phyxminiproblems-problemphyxmini0491-legprocesskind, def:physics:phyx-mini-0491:phyxminiproblems-problemphyxmini0491-idealgasbdaprocess, def:physics:phyx-mini-0491:phyxminiproblems-problemphyxmini0491-internalenergyinjoules, def:physics:phyx-mini-0491:phyxminiproblems-problemphyxmini0491-workdonebygasinjoules, def:physics:phyx-mini-0491:phyxminiproblems-problemphyxmini0491-heattransferredintogasinjoules}
  Macroscopic laws used by the calculation: ideal-gas internal energy depends only on absolute temperature; isobaric boundary work is p (V\_f - V\_i); isovolumetric boundary work vanishes; on each leg, Q\_in = U\_f - U\_i + W\_by. No field states a leg's numerical heat, the total heat, or the selected answer
\end{definition}
\begin{proof}
  This is the declared model definition of Satisfies Ideal Gas Process Laws.
\end{proof}
\begin{definition}[Answer Choice]
  \label{def:physics:phyx-mini-0491:phyxminiproblems-problemphyxmini0491-answerchoice}
  \lean{PhyXMiniProblems.ProblemPhyXMini0491.AnswerChoice}
  Labels of the four answers printed with the problem
\end{definition}
\begin{proof}
  This is the declared model definition of Answer Choice.
\end{proof}
\begin{definition}[displayed Total Heat In Joules]
  \label{def:physics:phyx-mini-0491:phyxminiproblems-problemphyxmini0491-displayedtotalheatinjoules}
  \lean{PhyXMiniProblems.ProblemPhyXMini0491.displayedTotalHeatInJoules}
  \uses{def:physics:phyx-mini-0491:phyxminiproblems-problemphyxmini0491-answerchoice}
  Joule value printed beside each answer label
\end{definition}
\begin{proof}
  This is the declared model definition of displayed Total Heat In Joules.
\end{proof}
\begin{definition}[recorded Answer Choice]
  \label{def:physics:phyx-mini-0491:phyxminiproblems-problemphyxmini0491-recordedanswerchoice}
  \lean{PhyXMiniProblems.ProblemPhyXMini0491.recordedAnswerChoice}
  \uses{def:physics:phyx-mini-0491:phyxminiproblems-problemphyxmini0491-answerchoice}
  Answer label recorded in the supplied dataset
\end{definition}
\begin{proof}
  This is the declared model definition of recorded Answer Choice.
\end{proof}
\begin{definition}[Rounds To Nearest Hundred Joules]
  \label{def:physics:phyx-mini-0491:phyxminiproblems-problemphyxmini0491-roundstonearesthundredjoules}
  \lean{PhyXMiniProblems.ProblemPhyXMini0491.RoundsToNearestHundredJoules}
  The source reports energy to the nearest hundred joules. This predicate says that the exact SI value lies within half of one 100 J display interval
\end{definition}
\begin{proof}
  This is the declared model definition of Rounds To Nearest Hundred Joules.
\end{proof}
\begin{definition}[Is Unique Closest Displayed Heat]
  \label{def:physics:phyx-mini-0491:phyxminiproblems-problemphyxmini0491-isuniqueclosestdisplayedheat}
  \lean{PhyXMiniProblems.ProblemPhyXMini0491.IsUniqueClosestDisplayedHeat}
  \uses{def:physics:phyx-mini-0491:phyxminiproblems-problemphyxmini0491-answerchoice, def:physics:phyx-mini-0491:phyxminiproblems-problemphyxmini0491-displayedtotalheatinjoules}
  A choice is uniquely closest to the exact derived heat
\end{definition}
\begin{proof}
  This is the declared model definition of Is Unique Closest Displayed Heat.
\end{proof}
\begin{lemma}[boundary Work Along BDA]
  \label{lem:physics:phyx-mini-0491:phyxminiproblems-problemphyxmini0491-boundaryworkalongbda}
  \lean{PhyXMiniProblems.ProblemPhyXMini0491.boundaryWorkAlongBDA}
  \uses{def:physics:phyx-mini-0491:phyxminiproblems-problemphyxmini0491-joulesperliteratmosphere, def:physics:phyx-mini-0491:phyxminiproblems-problemphyxmini0491-idealgasbdaprocess, def:physics:phyx-mini-0491:phyxminiproblems-problemphyxmini0491-workdonebygasinjoules, def:physics:phyx-mini-0491:phyxminiproblems-problemphyxmini0491-totalworkdonebygasinjoules, def:physics:phyx-mini-0491:phyxminiproblems-problemphyxmini0491-matchesproblemandprimaryfigure, def:physics:phyx-mini-0491:phyxminiproblems-problemphyxmini0491-satisfiesidealgasprocesslaws}
  The two boundary-work contributions are -16 L atm and zero. These values are derived conclusions, not fields of a premise structure
\end{lemma}
\begin{proof}
  The conclusion follows from the governing relations and figure readouts named in the statement of boundary Work Along BDA.
\end{proof}
\begin{lemma}[internal Energy Returns To Initial Value]
  \label{lem:physics:phyx-mini-0491:phyxminiproblems-problemphyxmini0491-internalenergyreturnstoinitialvalue}
  \lean{PhyXMiniProblems.ProblemPhyXMini0491.internalEnergyReturnsToInitialValue}
  \uses{def:physics:phyx-mini-0491:phyxminiproblems-problemphyxmini0491-idealgasbdaprocess, def:physics:phyx-mini-0491:phyxminiproblems-problemphyxmini0491-internalenergyinjoules, def:physics:phyx-mini-0491:phyxminiproblems-problemphyxmini0491-matchesproblemandprimaryfigure, def:physics:phyx-mini-0491:phyxminiproblems-problemphyxmini0491-satisfiesidealgasprocesslaws}
  Equal endpoint temperatures give zero total internal-energy change
\end{lemma}
\begin{proof}
  The conclusion follows from the governing relations and figure readouts named in the statement of internal Energy Returns To Initial Value.
\end{proof}
\begin{lemma}[total Heat Along BDA exact]
  \label{lem:physics:phyx-mini-0491:phyxminiproblems-problemphyxmini0491-totalheatalongbda-exact}
  \lean{PhyXMiniProblems.ProblemPhyXMini0491.totalHeatAlongBDA_exact}
  \uses{def:physics:phyx-mini-0491:phyxminiproblems-problemphyxmini0491-joulesperliteratmosphere, def:physics:phyx-mini-0491:phyxminiproblems-problemphyxmini0491-idealgasbdaprocess, def:physics:phyx-mini-0491:phyxminiproblems-problemphyxmini0491-totalheatintogasinjoules, def:physics:phyx-mini-0491:phyxminiproblems-problemphyxmini0491-matchesproblemandprimaryfigure, def:physics:phyx-mini-0491:phyxminiproblems-problemphyxmini0491-satisfiesidealgasprocesslaws}
  The first law and cancellation of the intermediate internal energy give an exact heat input of -16 L atm = -8106/5 J = -1621.2 J
\end{lemma}
\begin{proof}
  The conclusion follows from the governing relations and figure readouts named in the statement of total Heat Along BDA exact.
\end{proof}
% --- Archon declaration topology end ---
% --- Archon physics formalization source end ---
